% Lösungen Einheit 4 – Ausklammern (zusammenbau v1.2)
\kbabschnitt{Ausklammern}

\lbloes{16.}{a)}{die Zahl 4}
\lbloes{}{b)}{die Variable $y$}
\lbloes{}{c)}{die Zahl 2 und die Variable $a$, also $2a$}
\lbloes{}{d)}{die Zahl 3}
\lbloes{17.}{a)}{$4 \cdot 2x + 4 \cdot 5$}
\lbloes{}{b)}{$2 \cdot 3a + 2 \cdot 7$}
\lbloes{}{c)}{$5 \cdot 2y + 5 \cdot 3$}
\lbloes{}{d)}{$3 \cdot 3b + 3 \cdot 2$}
\lbloes{18.}{a)}{$3 \cdot (3x + 7)$}
\lbloes{}{b)}{$2 \cdot (2y + 5)$}
\lbloes{}{c)}{$5 \cdot (2a + 7)$}
\lbloes{}{d)}{$7 \cdot (2b + 3)$}
\lbloes{19.}{a)}{$5 \cdot (x + 2)$}
\lbloes{}{b)}{$5 \cdot (x + 5)$}
\lbloes{}{c)}{$x \cdot (x + 5)$}
\lbloes{}{d)}{$3x \cdot (2x + 5)$}
\lbloes{}{e)}{$6 \cdot (x + 1)$}
\lbloes{}{f)}{$3 \cdot (x^2 + 2x + 5)$}
\lbloes{}{g)}{$6 \cdot (x + 5)$}
\lbloes{}{h)}{Die 20 wurde nicht durch 5 geteilt, die Probe ergibt $5x + 100$}
\lbloes{}{i)}{$3a \cdot (a + 2b + 3)$}
\lbloes{20.}{a)}{Nicht stimmen können die Rechnungen 2 und 3. 2: Im Term steht ein Minus, in der Klammer ein Plus. 3: Aus zwei Gliedern werden in der Klammer wieder zwei, hier fehlt die Eins.}
\lbloes{}{b)}{(1) falsch, z. B. $3x + 5$: 3 und 5 haben keinen gemeinsamen Faktor größer als 1. (2) wahr, denn Ausklammern und Ausmultiplizieren sind nach dem Distributivgesetz Umkehrungen. (3) wahr, denn jedes Glied wird durch den Faktor geteilt und bleibt stehen, notfalls als Eins.}
\lbloes{}{c)}{Zwei Rechtecke nebeneinander, beide 6 cm hoch, eines $b$ cm breit, eines 3 cm breit; zusammen $6b + 18$ cm²}
\lbloes{}{d)}{$675$ €}
